\documentclass[11pt]{article}
\usepackage[margin=.80in]{geometry}
\usepackage{amsmath,amssymb,amsthm,booktabs}
\usepackage[hidelinks]{hyperref}
\newtheorem{theorem}{Theorem}
\newtheorem{lemma}{Lemma}
\newcommand{\tr}{\operatorname{tr}}
\newcommand{\Sym}{\operatorname{Sym}_0}
\newcommand{\Res}{\operatorname{Res}}
\newcommand{\I}{\mathbb 1}
\newcommand{\cA}{\mathcal A}
\title{Positive string poles and indefinite finite-transfer\\ graviton amplitudes on the full physical polarization space}
\author{}
\date{}
\begin{document}
\maketitle
\begin{abstract}
In ten-dimensional Type-II sphere-tree graviton scattering, we give a
center-of-mass transport for the full 35-dimensional real physical
polarization space, including the eight states omitted by the common-transverse
restriction. For a real probe carried with the same label on its incoming and
outgoing legs, the crossed projectile block reduces to the identity, a
rank-one probe projector and a double commutator. At fixed dimensionless
spacelike transfer $-1/2\le y<0$, every massive pole and every pole weight
after subtraction of its leading Regge asymptote is positive definite.
Nevertheless the exact finite-transfer centered quadratic coefficient after
Einstein subtraction is indefinite for every pure probe, and can be strictly
negative in every channel for a positive mixture of probes. We classify its
spectrum and find the sharp mixing threshold, rank five. The two sign
statements coexist because removal of the infinite leading Regge tower and
the massless exchange leaves a finite scalar term. No general nonforward
optical-theorem inequality, forward Wilson-coefficient sign, or loop
extension is assumed.
\end{abstract}

\section{Setup}
We work at sphere tree level in flat ten-dimensional spacetime with
mostly-plus signature. Write $x=a s$, $y=a t$, $z=a u$, $a=\alpha'/4$,
and $x+y+z=0$. Our dimensionless matrix amplitude and its Einstein part are
\begin{equation}
 \cA=\frac{\mathcal N}{xyz}\mathcal V,\qquad
 \cA_{\rm E}=\frac{\mathcal N}{xyz},\qquad
 \mathcal V=\prod_{\xi=x,y,z}\frac{\Gamma(1-\xi)}{\Gamma(1+\xi)}.
 \label{eq:amplitude}
\end{equation}
The physical amplitude equals $(8\pi G_{10}/a)\cA$ in the convention in
which the universal forward $t$-channel residue on unit projectile and
probe tensors is $-8\pi G_{10}s^2$ times the identity. The positive
prefactor follows from canonically normalized Einstein exchange, and
not from a conversion between coupling symbols in a polarization
calculation. In particular $a$ and the positive prefactor do not enter
any inertia count. We use the standard Type-II four-NS--NS tensor and
Virasoro--Shapiro form factor \cite{pasukonis,t8review}. Previous
gravitational dispersion theory already treats the string graviton pole
and Regge term \cite{tokuda,noumi,hamada}; the result here concerns the
complete tied physical polarization block and its sign structure, not
a new amplitude or a universal positivity bound.

For clarity we use all incoming momenta and COM energy one in computing
the tensor. The momentum span has time and two spatial coordinates,
in which
\begin{equation}
 \begin{split}
 k_1=(1,0,1),&\quad k_2=(1,0,-1),\\
 k_3=(-1,-r,-c),&\quad k_4=(-1,r,c),\qquad r^2+c^2=1 .
 \end{split}
\end{equation}
The common spacelike orthogonal complement $U$ has dimension seven.
On the four individual physical transverse spaces use the common $U$
and the four real unit lines
\begin{equation}
 \ell_1=(0,1,0),\quad \ell_2=(0,-1,0),\quad
 \ell_3=(0,c,-r),\quad \ell_4=(0,-c,r).
 \label{eq:transport}
\end{equation}
Thus the outgoing pair is attached by the same spatial rotation and
each vector label lives in an abstract $U\oplus\mathbb R$. This is not a
common eight-dimensional subspace of spacetime. The graviton labels
form the real Hilbert space $H=\Sym(\mathbb R^8)$, with
$\langle E,F\rangle=\tr(EF)$ and $\dim H=35$. The real probe label $P$ is
unit and the same on legs 2 and 4; the external projectile matrix
pairs arbitrary labels on legs 1 and 3. Complex projectiles use the
sesquilinear extension, with the conjugate outgoing coefficients.
No claim is made for independently retuned probes on each leg or an
arbitrary two-particle crossing matrix.

\section{Full transverse transport and the operator numerator}
Use $F_i=(k_i e_i^{\mathsf T}-e_i k_i^{\mathsf T})g$ as a mixed
Lorentz-index field-strength matrix. For vector labels, the
normalized chiral numerator is $-t_8(F_1,F_2,F_3,F_4)/4$ in the
normalization of \eqref{eq:amplitude}. Here $t_8$ is eight times
the sum of the three single traces with cyclic representatives
$(1234),(1324),(1342)$, minus twice the sum of three double traces
corresponding to the pairings $(12)(34),(13)(24),(14)(23)$. Overall
normalization is fixed by the universal Einstein residue above.
This contraction is gauge-safe before substitution: replacing any
$e_i$ by $k_i$ gives zero. In the frames \eqref{eq:transport} it is
\begin{equation}
 yz (e_1\cdot e_2)(e_3\cdot e_4)
 +xz (e_1\cdot e_3)(e_2\cdot e_4)
 +xy (e_1\cdot e_4)(e_2\cdot e_3),
 \label{eq:vector}
\end{equation}
where the dots are in the abstract label space.
We give a finite verification of this identity. For zero planar
labels, the three common-transverse pairing coefficients are
$yz,xz,xy$. For one or three planar labels, reflection of $U$ changes
the sign and the expression is zero. With two planar labels at legs
$12,34$ the coefficient is $yz$; at $13,24$ it is $xz$; at $14,23$ it
is $xy$. In each case the remaining labels contract by the Euclidean
dot. With four planar labels the value is $yz+xz+xy$. These entries
follow by substituting \eqref{eq:transport} into the mixed-index trace
formula; for example $x=4,y=-2(1-c),z=-2(1+c)$ in the unit COM frame,
so the three two-planar coefficients are
$4(1-c^2),-8(1+c),8(c-1)$.
For zero planar labels, $O(7)$ invariance leaves the three Euclidean
pairings and two distinct common coordinate axes distinguish all
coefficients. For two planar labels only the common dot is possible.
Thus the elementary checks using two axes and one planar label exhaust
the invariant tensor types, rather than sampling the full state space.
The accompanying exact replay substitutes all $3^4$ corresponding
multilinear cases, reducing polynomials modulo $r^2+c^2-1$.
Homogeneity restores general scale. The polynomial
\eqref{eq:vector} gives the analytic continuation with these labels
held fixed; exchanging the same probe label implements $x\leftrightarrow z$
in this block before any dispersion argument.

For the closed tensor, apply \eqref{eq:vector} to left and right
labels and contract with the symmetric traceless leg tensors. Write
$Q=|P\rangle\langle P|$, $\mathcal G(E)=[P,[P,E]]$.
For example, the diagonal pairing terms across both copies are
$\tr(EP)\tr(FP),\tr(EF),\tr(EP)\tr(FP)$; the six crossed terms
are four instances of $\tr(EF P^2)$ and two of $\tr(EPFP)$.
Their coefficients sum to
\[
 \mathcal N=(xz)^2\I+y^2(x^2+z^2)Q-xzy^2\mathcal G.
 \label{eq:numer}
\tag{4}
\]
Indeed the two four-trace coefficients are $2xyz(x+z)$ and
$2xy^2z$; using $x+z=-y$ and
$\langle E,\mathcal G F\rangle=2\tr(EF P^2)-2\tr(EPFP)$ gives (4).
The matrix numerator is on shell. If $p_i$ are the eigenvalues of $P$,
then $\mathcal G$ vanishes on the commuting blocks and has
eigenvalues $(p_i-p_j)^2$ between eigenspaces. In particular
\begin{equation}
 0\preceq \mathcal G\preceq 2\I,\qquad Q\mathcal G=0,\qquad
 Q+\mathcal G/2\preceq\I,
 \label{eq:geom}
\end{equation}
since $2-(p_i-p_j)^2=(p_i+p_j)^2+
2\sum_{k\ne i,j}p_k^2$.

\section{Positive residues, including the Regge remainder}
Put $v=x+y/2$. At $x=n\ge1$, $\delta=-y\in(0,1/2]$, recurrence
for the gamma factors gives
\[
 -\Res_{x=n}\cA
 =R_n(y)\,\mathcal N\big|_{x=n},\qquad
 R_n(y)=\left[\frac{\Gamma(n+y)}
                  {\Gamma(1+y)\Gamma(n+1)}\right]^2.
\]
The positive-$x$ Feynman prescription therefore has $+\pi$ times
this matrix as the pole spectral coefficient. The crossed pole pair
contributes $W_n=2R_n\mathcal N|_{x=n}/(n+y/2)^3$ to the centered
twice-subtracted coefficient after the two massless $s,u$ poles
are removed. Explicitly $W_n=w_n\I+\alpha_n Q+\beta_n\mathcal G$ with
\begin{equation}
 \begin{split}
 w_n&=\frac{2R_n n^2(n+y)^2}{(n+y/2)^3},\\
 \alpha_n&=\frac{2R_n y^2[n^2+(n+y)^2]}{(n+y/2)^3},\\
 \beta_n&=\frac{2R_n y^2 n(n+y)}{(n+y/2)^3}.
 \end{split}
 \label{eq:weights}
\end{equation}
All are strictly positive. This is a property of the specific
amplitude in this band, not a nonforward optical theorem.

The leading large-$n$ term follows from the gamma ratio and (4):
$L_n=\ell_n\I$, $\ell_n=2 n^{2y-1}/\Gamma(1+y)^2$.
Let $D_n=W_n-L_n$. We prove $w_n>\ell_n$ independently of
polarization. The ratio starts at one at $\delta=0$, and
\begin{equation}
 \partial_\delta\log(w_n/\ell_n)
 =-\frac2{n-\delta}+\frac3{2n-\delta}
      +2[\log n-\psi(n-\delta)]
 > \frac{n-2\delta}{(n-\delta)(2n-\delta)}
      -2\log(1-\delta/n)>0 .
\end{equation}
Here $\psi(u)<\log u-1/(2u)$ for $u>0$ by its standard positive
integral representation \cite{nist}; this applies at the endpoint
$n=1,\delta=1/2$ as well. Hence $D_n\succ0$.
Moreover $\alpha_n-2\beta_n=2R_ny^4/(n+y/2)^3>0$.
By \eqref{eq:geom} the sharp maximum of every $D_n$ is
$w_n-\ell_n+\alpha_n$, attained uniquely on $P$ on the full space;
the minimum is $w_n-\ell_n$. The same maximal direction holds for any
tail sum with this pure probe. At fixed $y<0$ both series of
$W_n$ and $L_n$ converge by their gamma asymptotics; the remainder also
has scalar leading-size bound via the sharp maximum. We do not require
a finite-level sign claim for the final subtracted contact.

\section{A finite scalar term separates the two sign questions}
The exact massless $s,u$ residues are $-y^2Q$, and the $t$-channel
residue is $-x^2\I$. Define
\begin{equation}
 \cA_{su}=\cA+y^2Q(1/x+1/z),\qquad
 C_P(y)=[v^2]\big(\cA-\cA_{\rm E}\big)
       =[v^2]\cA_{su}+ \frac1y\I.
 \label{eq:contactdef}
\end{equation}
The quantity $C_P$ is a finite-transfer centered analytic coefficient,
not an isolated local EFT Wilson coefficient. At fixed $y<0$, applying
Cauchy's theorem to $\cA_{su}/v^3$ and pairing crossed poles gives
$[v^2]\cA_{su}=\sum_{n\ge1}W_n$. The large contour contribution vanishes:
away from the real pole sequences and on contours crossing at a fixed
distance from those poles, gamma reflection and ratio asymptotics
bound $\cA/v^2$ by $O(|v|^{2y})$. The additional massless subtraction
decays on the same contours. In particular the massive sum is taken
before the forward limit. Therefore
\begin{equation}
 C_P(y)=\left(\frac{2\zeta(1-2y)}{\Gamma(1+y)^2}+\frac1y\right)\I
            +\sum_{n\ge1}D_n(y).
 \label{eq:counter}
\end{equation}
The coefficient of $\I$ is a finite term at fixed nonzero transfer.
In particular it is not legitimate to discard it when drawing sign
conclusions from the positive matrices $D_n$. Equation
\eqref{eq:counter} sums a convergent positive tower at negative
transfer and does not use zeta regularization.

The exact sign analysis can also be derived directly without this
contour. Define
\begin{equation}
 \begin{split}
 V_0(y)&=\frac{\Gamma(1-y)}{\Gamma(1+y)}
       \left[\frac{\Gamma(1+y/2)}{\Gamma(1-y/2)}\right]^2,\qquad
 H_1(y)=\psi_1(1+y/2)-\psi_1(1-y/2),\\
 U&=-\log V_0(-\delta),\qquad
 A=\frac{1-e^{-U}}{\delta},\qquad b=\delta e^{-U} H_1(-\delta)>0 .
 \end{split}
 \label{eq:Ab}
\end{equation}
Here $\psi_1$ denotes the trigamma function, and
$\mathcal V=V_0(y)(1+H_1(y)v^2+O(v^4))$.
At the center $\cA_{\rm E}(0,y)=y(\I/4+2Q-\mathcal G)$ and
$[v^2]\cA_{\rm E}=(-\I+16Q)/y$ by (4).
Insert these into $[v^2](\mathcal V-1)\cA_{\rm E}$ to find
\begin{equation}
 C_P=c_0\I+(16A-2b)Q+b\mathcal G,\qquad c_0=-A-b/4 .
 \label{eq:contact}
\end{equation}

\begin{theorem}[Full physical inertia]
If the eigenspace dimensions of a unit real $P$ are $m_i$ with
distinct eigenvalues $p_i$, the exact contact eigenvalues are
\begin{center}
\begin{tabular}{lll}
\toprule
Sector & Multiplicity & Value\\
\midrule
$P$ & $1$ & $15A-9b/4$\\
Commutes with $P$, orthogonal to $P$ &
 $\sum_i m_i(m_i+1)/2-2$ & $c_0$\\
Between eigenspaces $i<j$ & $m_i m_j$ &
 $b[(p_i-p_j)^2-\theta]$, \ $\theta=1/4+A/b$\\
\bottomrule
\end{tabular}
\end{center}
Every pure probe gives both signs. A probe with four eigenvalues
$1/\sqrt8$ and four $-1/\sqrt8$ gives exactly one positive and 34 negative
channels. In particular the two-positive floor on the narrower
seven-axis common-transverse problem does not extend to this
physical state space.
\end{theorem}
\begin{proof}
Use the block decomposition of $\mathcal G$ above and that its kernel
contains $P$. The two strict pure-probe signs follow from $A/b>1/4$
proved below and the commuting complement of dimension at least six.
For the stated probe every distinct squared gap is $1/2<\theta$.
\end{proof}

\section{Mixed probes and the sharp threshold}
Let $\rho=\sum_jw_j|P_j\rangle\langle P_j|$ be a real positive trace-one
density on $H$, implemented by using identical $P_j$ labels on the two
probe legs at each member of the ensemble and averaging the resulting
projectile block. Denote $\mathcal G_\rho=\sum_jw_j\mathcal G_{P_j}$.
Equation \eqref{eq:geom} gives $\rho+\mathcal G_\rho/2\preceq\I$;
in particular residue and subtracted-residue positivity persist.
On the other hand
\begin{equation}
 C_\rho=c_0\I+(16A-2b)\rho+b\mathcal G_\rho,\qquad
 \tr C_\rho=-19A-11b/4<0.
\end{equation}
The trace follows from $\tr_H\mathcal G_P=\sum_{i<j}(p_i-p_j)^2=8$.
Thus no positive probe ensemble makes the contact positive semidefinite.

\begin{theorem}[Sharp real mixing rank]
For $1\le r\le 7$ the exact minimax is
\[
 \min_{\rho\succeq0,\ \tr\rho=1,\ \operatorname{rank}\rho\le r}
       \lambda_{\max}(C_\rho)=c_0+\frac{16A-2b}{r}.
\]
The minimizing densities are exactly the normalized projections onto
$r$-dimensional mutually commuting real symmetric subspaces of $H$.
Rank five is the least rank for which all 35 contact channels can
become strictly negative, at every fixed transfer in the band.
\end{theorem}
\begin{proof}
Positivity of $\mathcal G_\rho$ and $\lambda_{\max}(\rho)\ge1/r$
give the lower bound. An equal mixture on an $r$-dimensional
traceless diagonal subspace attains it: the commuting sector is
unchanged by $\mathcal G_\rho$, while
$\mathcal G_\rho$ on the $(i,j)$ off-diagonal channel is
$\|\Pi(e_i-e_j)\|^2/r\le2/r$, for $\Pi$ the corresponding diagonal
projection. Since $16A-2b>2b$, this cannot surpass the top support
value. If the lower bound is attained, $\rho$ must be the normalized
projection of rank $r$. On that support each quadratic form
$\sum_jw_j\|[P_j,E]\|_{\rm F}^2$ must vanish, proving the commuting
condition. Conversely such a subspace diagonalizes simultaneously.

For the signs, one needs precisely the strict uniform scalar estimate
\begin{equation}
  \frac14<\frac{A}{b}<\frac{13}{44}.
 \label{eq:ratio}
\end{equation}
Indeed the upper bound gives strict negativity at rank five and
the lower bound gives positivity at rank four. To prove \eqref{eq:ratio}, use
the convergent series from \eqref{eq:Ab}:
\[
 U=\sum_{k\ge1}\frac{2(1-4^{-k})\zeta(2k+1)}{2k+1}\delta^{2k+1},
 \qquad
 \delta^2H_1(-\delta)=\sum_{k\ge1}8k4^{-k}\zeta(2k+1)\delta^{2k+1}.
\]
Their coefficient ratios are
$(4^k-1)/(4k(2k+1))$, equal to $1/4$ at $k=1$ and
strictly higher thereafter. Since $(e^U-1)/U>1$ the lower bound follows.
For the upper bound put $t=\delta^2\le1/4$ and
$w=t\zeta(5)/\zeta(3)\le1/4$. Bounding coefficients of the
numerator after $k=2$ by $2\zeta(3)/7$ gives
\[
 \frac{U}{\delta^2 H_1}
 \le \frac{1/2+3w/8+(2/7)t^2/(1-t)}{2+w}
 \le \frac{1/2+3/32+1/42}{9/4}.
\]
The ratio increases in $w$ throughout this rectangle, including
after maximizing the displayed tail. Separately
$U\le\zeta(3)\delta^3/2+8\zeta(5)\delta^5/15<153/1600$
by bounding the geometric tail and using
$\zeta(3)<5/4,\zeta(5)<21/20$.
Thus $(e^U-1)/U\le (1-U/2)^{-1}<3200/3047$ by the
factorial series. The product bounds $A/b$ by
$166000/575883<13/44$ as required.
\end{proof}

\section{Interpretation and limitations}
The separation of signs uses the physical projected amplitude with
a common real probe label, not a sum of positive on-shell probabilities.
The analytically prescribed massless and Regge subtractions must be
retained together. The actual sign-changing contact operator and the
positive spectral remainder are both determined by the same full
little-group data. The wider state space is therefore substantive:
it changes the pure-probe inertia floor and supplies a complete
polarization description of the cancellation mechanism.

The transfer $y=0$ itself is not part of the fixed-transfer contour
argument. In particular negative finite-transfer coefficients neither
violate tree unitarity nor imply a negative forward coefficient.
Likewise no universal sign for amplitudes on independent two-particle
polarizations, a compactified dimension, or an interacting loop
amplitude is inferred. The inverse scale-identification problem
studied in a projected Type-II block \cite{scale} uses different
observables and does not establish this polarization theorem.

\begin{thebibliography}{9}
\bibitem{pasukonis} J.~Pa\v{s}ukonis, \emph{Gravitational scattering of massless scalars in QFT and superstring theory}, \href{https://arxiv.org/abs/hep-th/0506065}{hep-th/0506065},
Secs.~2--3.
\bibitem{t8review} C.~R.~Mafra and O.~Schlotterer,
\emph{Tree-level amplitudes from the pure spinor superstring},
\href{https://arxiv.org/abs/2210.14241}{arXiv:2210.14241}, Sec.~8.
\bibitem{tokuda} J.~Tokuda, K.~Aoki and S.~Hirano,
\emph{Gravitational positivity bounds},
\href{https://arxiv.org/abs/2007.15009}{arXiv:2007.15009}, Sec.~3.3.
\bibitem{noumi} T.~Noumi and J.~Tokuda,
\emph{Finite energy sum rules for gravitational Regge amplitudes},
\href{https://arxiv.org/abs/2212.08001}{arXiv:2212.08001}.
\bibitem{hamada} Y.~Hamada, R.~Kuramochi, G.~J.~Loges and S.~Nakajima,
\emph{On (scalar QED) gravitational positivity bounds},
\href{https://doi.org/10.1007/JHEP05(2023)076}{JHEP 05 (2023) 076}.
\bibitem{nist} NIST Digital Library of Mathematical Functions,
\href{https://dlmf.nist.gov/5.9.E15}{Eq.~5.9.15}.
\bibitem{scale} \emph{Recovering Type-II string scales from nonforward
five-dimensional graviton scattering},
\href{https://github.com/g3173028019-beep/ai-automated-quantum-gravity-research/tree/main/papers/32-typeii-scale-recovery}{research archive, Paper 32}.
\end{thebibliography}
\end{document}
